%%%%%%%%%%%%%%%%%%%%%%%%%%%%%%%%%%%%%%%%%%%%%%%%%%%%%%%%%%%%%%%%%%%%%%
%  第 1 个逻辑月 · 月度报告
%  主题：小块：从抽象到具体（复现学长 + UDF 基线 + 算法提要）
%
%  说明：本报告由三人的既有材料按组长提纲的逻辑月重排而成。
%        尚无素材之处按规范留空并标注，不编造。
%  编译：xelatex 本文件（跑两遍）
%%%%%%%%%%%%%%%%%%%%%%%%%%%%%%%%%%%%%%%%%%%%%%%%%%%%%%%%%%%%%%%%%%%%%%
\documentclass[stage1,fontset=fandol]{ydzjreport}

\stageonesetup
{第 1 个逻辑月 · 基础建立}%
{冯旭东、陈思文}%
{复合微通道分区冷却}%
{【待确认】逻辑月 1 对应的自然月}%
{小块划分；学长复现；水物性 UDF；算法提要}%
{确立小块的划分方式，复现学长的建模与物理性质（UDF）思路，
	并完成控制方向的算法提要。
}%

\begin{document}
	\makestageonehead
	
	% ==================================================================
	\section{本月研究任务}
	
	\subsection{任务目标}
	
	本月项目的核心任务关键词为
	\textbf{「复现已有调研论文+课题组学长提供资料的结构、了解液体换热物理性质，编写对应变物性UDF，初步搭建冷板思路，初步构建控制流程」}、
	据此分解为四项任务：
	
	\begin{enumerate}
		\item 通过阅读文献确立\textbf{冷板的优化方式}，把空间不均匀热需求落实为研究的结构；
		\item 复现已有文献、课题组学长在\textbf{结构建模}上的思路；
		\item 复现已有文献、课题组学长在\textbf{物理性质（UDF）}上的设置，并核查其正确性；
		\item 初步控制方向完成\textbf{流程框架构建}。
	\end{enumerate}
	
	\subsection{完成情况}
	
	\begin{table}[H]
		\centering
		\caption{本月任务完成情况}
		\label{tab:tasks}
		\begin{tabularx}{\textwidth}{@{}Y c Y@{}}
			\toprule
			\rowcolor{ydzjlight} 任务 & 状态 & 说明 \\
			\midrule
			小块划分方式确立 & 部分完成 & 已形成概念表述，具体结构待进一步探究 \\
			学长结构建模复现 & 已完成 & 已复现论文中的多种流道布局 \\
			学长物性（UDF）思路复现 & 已完成 & 核查后发现原函数存在误差 \\
			控制算法提要 & 已完成 & 文献要点与二阶对象特性已整理 \\
			\bottomrule
		\end{tabularx}
	\end{table}
	
	
	% ==================================================================
	\section{核心概念梳理}
	
	\subsection{概念一：非均匀热载荷与“拼图式”小块划分}
	
	芯片热流密度在空间上高度不均匀：结区集中发热，背景区几乎不发热。
	若采用均匀结构，背景区会被「一视同仁」地强力冷却，形成\textbf{过度冷却}——
	既拉大温度梯度，又浪费泵功。
	
	\textbf{“拼图式”小块}的作用，是把连续冷板上的空间不均匀热需求，转化为
	\textbf{可单独研究、可单独供液}的流道单元。
	
	需要说明的是：各小块在连续冷板上形成不同几何区域，
	\textbf{并不必然意味着把冷板机械切开或完全隔热}。
	先研究单元、再研究组合，有利于识别各形态的作用，
	并为独立供液或流量调控预留条件。
	
	\textcolor{red}{【待补：小块划分概念图示意图（王佳蔚）】}
	
	\subsection{概念二：物理性质是探究的起点}
	
	（1）冷却工质水的物性随温度变化。因此\textbf{物性核查必须放在结构比较之前}。\textcolor{red}{补充UDF的编写学习过程，补充在后面，这里不需要（陈思文、王佳蔚）}
	（2）控制策略的物理起源\textcolor{red}{补充控制的缘由和基础（张昊洋）、控制框图（王佳蔚）}
	\textcolor{red}{补充初步控制框架}
	
	% ==================================================================
	\section{阶段一：通读与摘录}
	
	\subsection{学习路径}
	
	\begin{enumerate}
		\item 先通读学长攻略（课题组已有经验）、结构类论文，梳理几何布局、进出口组织与边界条件；
		\item 复现其中的代表性流道，取得可对照的性能趋势；
		\item 转向物理性质：核查计算所用的水物性函数，确认随温度变化的正确性；
		\item 控制方向并线推进：整理控制文献要点与二阶对象的动态特性，形成算法提要。
	\end{enumerate}
	
	\subsection{关键摘录}
	
	

	\begin{table}[H]
	  \centering
	  \small
	  \renewcommand{\arraystretch}{1.12}
	  \caption{文献要点与本阶段工作}
	  \label{tab:notes}
	  \begin{tabularx}{\textwidth}{@{}>{\RaggedRight\arraybackslash}p{0.16\textwidth} Y Y@{}}
	    \toprule
	    \rowcolor{ydzjlight} 文献 & 论文要点 & 本阶段的应用 \\
	    \midrule
	    Ao 等\cite{ao} & 方形、圆形与六边形微通道的比较 & 复现三种流道，学习短流路与局部强化的组织方式 \\[3pt]
	    隋增光等\cite{ref1} & 综述数据中心吸收式器件散热、余热制冷、储能、发电与空气取水，强调低品位余热的多功能利用及系统集成。 & 补充研究背景与系统应用场景；在冷板研究之外，梳理出口水温、余热利用与节能节水之间的联系，拓展系统评价思路。 \\[3pt]
	    Xu 等\cite{ref2} & 针对双桥臂 SiC IGBT 多热源，研究双面液冷与六类导流微结构；海鸥形结构综合性能因子达 1.51，交错结构更利于芯片局部均温。 & 依据热源位置布置强化结构，将海鸥形、交错导流纳入候选矩阵；比较模块整体温度与局部热点温度，避免仅以平均温度筛选。 \\[3pt]

	    \bottomrule
	  \end{tabularx}
	\end{table}

	\begin{table}[H]
	  \centering
	  \small
	  \renewcommand{\arraystretch}{1.12}
	  \caption{文献要点与本阶段工作（续1）}
	  \label{tab:notes-cont-2}
	  \begin{tabularx}{\textwidth}{@{}>{\RaggedRight\arraybackslash}p{0.16\textwidth} Y Y@{}}
	    \toprule
	    \rowcolor{ydzjlight} 文献 & 论文要点 & 本阶段的应用 \\
	    \midrule
	    Wang 等\cite{ref5} & 将非等高双层通道、分形分叉与冲击射流结合，并采用实验验证和 NSGA-II 优化；文中报告最优无肋方案的 PEC 较参照结构提高 24.61\%。 & 拓展“层间高度差—分叉配流—局部冲击”的组合方案；先比较不同进水组织，再开展热性能与压降的双目标参数优化。 \\[3pt]
	    张乾坤等\cite{ref6} & 比较普通二级分支与三类叶脉仿生流道；完整五级分支相对普通结构压降降低超过 35\%，传热面温度降低超过 2 $\,{}^\circ\mathrm{C}$，综合强化因子大于 1.2。 & 将分支级数、分叉形式与网络完整性加入结构矩阵；比较流量分配、温度均匀性和压降，探索以配流优化实现强化与降阻。 \\[3pt]
	    Yang 等\cite{ref7} & 综述增材制造微通道散热器的材料、工艺、复杂流道及优化方法，讨论分布式歧管、多层结构、仿生网络与制造约束之间的关系。 & 建立候选结构的可制造性记录，关注最小特征尺寸、内壁粗糙度、残粉清理与密封；筛选既有热水力优势又便于加工验证的方案。 \\[3pt]
	    Hu 等\cite{ref8} & 比较三类吸液芯的异形均热板；顺序梯度网芯兼顾毛细回流与渗透性，最高测试热流密度达 200 W/cm$^2$，对应热阻为 0.047 K/W。 & 借鉴梯度分区与高低阻通路协同设计；将均热板作为复合散热的后续备选，并区分其相变传热机理与本阶段单相微通道模型。 \\[3pt]
	   
	    
	    Ansari 与 Kim\cite{ref13} & 针对非均匀热流，背景区使用矩形微通道、热点区布置圆柱针肋；以热阻、泵功、温度非均匀性和热点最大温升评价分区冷却。 & 建立背景区与热点区分开的热源边界条件；将局部针肋强化与全域强化进行对照，考察热点抑制效果及其额外泵功代价。 \\[3pt]
	    
	    Ma 等\cite{ref15} & 实验与模拟比较等周期和变周期正弦微通道；变周期结构增强扰动并改善均温性，但压降增大；在相同泵功下，平均热阻降低约 13\%。 & 将周期沿程变化加入几何参数矩阵；同时开展相同流量与相同泵功比较，分析下游强化能否补偿冷却液温升带来的散热衰减。 \\[3pt]
	    
	    王雅博等\cite{ref16} & 比较两端进出水与中心垂直射流配合一、二、四出口的布局；中心射流四出口方案虽非 Nu 最高，但流阻较小、综合性能更优且可避免局部热点。 & 将进口位置与出口数量作为独立变量；固定内部结构比较不同配流方式，联合考察速度死区、温度均匀性、Nu 和压降。 \\[3pt]
	    
	    Choi 等\cite{ref17} & 通过实验与 58 组数值算例研究大面积风冷 N 型歧管微通道；分析通道宽、长、高对流量分配、热阻与压降的影响，突出短流路和一体化结构的作用。 & 借鉴 N 型歧管分流与短流路设计，补充通道尺寸对配流的敏感性分析；应用于液冷时重新计算工质相关结果，保留其拓扑设计思路。 \\[3pt]
	    \bottomrule
	  \end{tabularx}
	\end{table}



	\begin{table}[H]
	  \centering
	  \small
	  \renewcommand{\arraystretch}{1.12}
	  \caption{文献要点与本阶段工作（续2）}
	  \label{tab:notes-cont-5}
	  \begin{tabularx}{\textwidth}{@{}>{\RaggedRight\arraybackslash}p{0.16\textwidth} Y Y@{}}
	    \toprule
	    \rowcolor{ydzjlight} 文献 & 论文要点 & 本阶段的应用 \\
	    \midrule
	    Indhe 等\cite{ref20} & 综述 CFD、分子动力学和人工智能在微电子散热中的应用，讨论微通道、纳米增强相变材料及复合冷却的多尺度分析与优化方法。 & 形成“参数化建模—CFD 数据积累—代理模型—优化验证”的后续路线；本阶段优先统一算例与数据字段，为机器学习准备可追溯样本。 \\[3pt]
	    Nie 等\cite{ref21} & 通过针肋关联导流结构削弱尾部流动死区并促进混合，但存在明显压降代价；在 0.15—0.8 mW 相同泵功范围内，最优结构热阻较传统针肋降低 20.8\%—22.2\%。 & 在针肋后方引入导流构件，考察高度、直径及相对偏移；结合速度场辨识强化机理，并以相同泵功对照判断额外阻力是否值得。 \\[3pt]
	    Yin 等\cite{yin} & 提出同轴开口环形 TSV 针肋，比较直列、错列与传统圆柱 TSV；综合分析 Nu、最高温度、压降、PEC、循环热应力及信号完整性。 & 借鉴开口环形针肋与错列组织方式；统一提取平均 Nu、最高温度、压降和 PEC，并将热应力作为封装可靠性的后续扩展指标。 \\[3pt]
	    Seong 等\cite{seong} & 研究倾斜肋—针肋的直列、错列及倾角效应；Re=450、倾角 60$^\circ$、热流密度 100 W/cm$^2$ 时，PEC 为 1.52，并建立 ANN 性能预测模型。 & 建立“排列方式$\times$倾角$\times$Re”参数矩阵，分析层间混合与停滞区变化；积累 CFD 数据后训练代理模型，并用独立算例验证预测结果。 \\[3pt]
	    Lorenzini 等\cite{ref24} & 在 200 $\mu$m 硅微间隙中按热点分布局部加密针肋；背景热流密度 250 W/cm$^2$、热点 250—750 W/cm$^2$、入口水温 21.3 $\,{}^\circ\mathrm{C}$ 时，所研究器件最高温度低于 65 $\,{}^\circ\mathrm{C}$。 & 建立均匀针肋与热点局部加密的对照模型；针对不同热点位置和强度匹配针肋密度，并核查旁路流动是否削弱热点区冷却。 \\[3pt]
	   
	    \bottomrule
	  \end{tabularx}
	\end{table}



	% ==================================================================
	\section{阶段二：导出与复现}
	
	\subsection{复现目标}
	
	本月复现两项内容：
	\textbf{(1)} 学长论文中的方形、圆形与六边形三种微通道的结构与流动换热组织；
	\textbf{(2)} 计算所用的水物性函数（UDF），核查其随温度变化的正确性。
	
	\subsection{复现环境与步骤}
	
	
	
	\textcolor{red}{补充复现图片和后处理图片（冯旭东），可以导入学长攻略参考}
	
	\textbf{环境}：ANSYS Fluent，工质为水，固体域与流体域耦合求解；
	物性核查采用 Python 后处理。
	
	\textbf{步骤}：
	
	\begin{enumerate}
		\item 按论文的方形、圆形与六边形布局分别构造流道，梳理流体域、固体域与加热区域；
		\item 对照论文的边界条件与求解方法开展计算，检查残差与监测量；
		\item 对三种结构做趋势对照，关注热负荷与流量变化下的温度响应；
		\item 核查原水物性 UDF：在 298.15 K 校验点比对参考物性值；
		\item 记录误差量级，判断是否需要重新拟合。
	\end{enumerate}
	
	\subsection{复现结果}
	
	\textbf{结构复现}：三种流道的复现已完成，计算结果与论文的趋势一致——
	流量增加时结温总体下降、压降增加；热流密度增加时结温上升。
	
	\begin{table}[H]
		\centering
		\caption{三种流道的复现情况}
		\label{tab:repro}
		\begin{tabularx}{\textwidth}{@{}Y Y Y@{}}
			\toprule
			\rowcolor{ydzjlight} 流道形态 & 复现关注内容 & 结果 \\
			\midrule
			方形微通道 & 方形流路及中心扰动结构的影响 & 已完成，趋势与论文一致 \\
			圆形微通道 & 环形流路及流体分配的影响 & 已完成，趋势与论文一致 \\
			六边形微通道 & 六边形布局及局部流动组织 & 已完成，趋势与论文一致 \\
			\bottomrule
		\end{tabularx}
	\end{table}
	
	
	
	\textbf{物性核查}：在 298.15 K 校验点，原水物性函数相对参考值
	\textbf{密度偏低约 60.8\%、比热偏低约 10.5\%、黏度偏低约 26.7\%}。
	
	这是一个\textbf{必须记录的重要发现}：它意味着此前所有基于该物性的
	温度与压差结果都存在系统偏差，结构之间的排序也可能受影响。
	本月据此完成了新物性 UDF 的编写与编译，挂接与复算列入后续月份。
	
	【缺数据：新 UDF 挂接后的物性自检结果】\textcolor{red}{王佳蔚补充新的的说明}
	
	% ==================================================================
	\section{阶段性收获}
	
	\subsection{认知上的变化}
	
	\textbf{一、从「整体冷却」转向「按需分区」。}
	原来倾向于把冷板当作一个整体来强化，认为加密结构总归有利。
	现在认识到：热流密度本身不均匀，均匀强化会在低负荷区造成过度冷却，
	\textbf{冷却资源必须按区域配置}。
	
	\textbf{二、物性不是背景条件，而是结论的前提。}
	复核原物性函数之前，我们默认它是正确的。
	发现偏差后才知道：\textbf{在物性不可信的情况下比较结构，等于在错的基准上排序}。
	因此把「物理性质」放在结构比较之前的做法是必要的。
	
	
	
	% ==================================================================
	\section{下月计划}
	
	\begin{table}[H]
		\centering
		\caption{下月工作计划}
		\label{tab:plan}
		\begin{tabularx}{\textwidth}{@{}Y Y c@{}}
			\toprule
			\rowcolor{ydzjlight} 事项 & 预期产出 & 时间 \\
			\midrule
			挂接新水物性 UDF 并完成物性自检 & 物性校验记录 & 第 1 周 \\
			补齐小块划分图，明确划分依据与数量 & 小块划分方案 & 第 1--2 周 \\
			建立统一的评价指标体系 & 指标定义表 & 第 2--3 周 \\
			小块内的形态与扰流单元仿真 & 形态对比数据 & 第 3--4 周 \\
			\textcolor{red}{补充自控学习}\\
		\end{tabularx}
	\end{table}
	
	% ==================================================================
	\begin{thebibliography}{99}
	  \bibitem{ao} Ao C, Xu B, Chen Z. Design and performance optimization of novel microchannel heat sinks for high-power electronic chips[J]. International Communications in Heat and Mass Transfer, 2026, 178: 111803.
	  \bibitem{ref1} 隋增光, 林昊晟, 隋云任, 等. 数据中心吸收式节能节水技术研究与应用进展[J]. 制冷学报, 2026, 47(1): 37-50. 
	  \bibitem{ref2} Xu H, Wang Y, Tian W, et al. Thermal performance optimization of dual-bridge arm SiC IGBT modules via microstructure-modified microchannel heat sinks under double-sided liquid cooling[J]. International Communications in Heat and Mass Transfer, 2026, 176: 111339. 
	  
	  
	  \bibitem{ref5} Wang M, Xie J, Bu J, et al. Thermo-hydraulic effects of asymmetric-height and fractal bifurcation channels combined with impingement in microchannel heat sinks[J]. Applied Thermal Engineering, 2026, 298: 130998. 
	  \bibitem{ref6} 张乾坤, 李红艳, 吕碧纯, 等. 叶脉仿生微通道换热器流动与传热特性的数值研究[J]. 制冷学报, 2025, 46(4): 52-60. 
	  
	  \bibitem{ref7} Yang G, Tang Y, Sun Y, et al. Fabrication of microchannel heat sinks using additive manufacturing: A review[J]. Renewable and Sustainable Energy Reviews, 2026, 233: 116879.
	  
	  \bibitem{ref8} Hu Z, Yan W, Wang S. Study on the Thermal Performance of a Specially-Shaped Vapor Chamber for Cooling High Heat Flux Chips[J/OL]. Thermochimica Acta, 2026.
	  

	  \bibitem{ref13} Ansari D, Kim K Y. Hotspot thermal management using a microchannel-pinfin hybrid heat sink[J]. International Journal of Thermal Sciences, 2018, 134: 27-39. 
	  
	  \bibitem{ref15} Ma D, Liu Y, Zhang X, et al. Experimental and numerical investigation of the thermohydraulic performance of a variable-period sinusoidal microchannel heat sink[J]. Applied Thermal Engineering, 2025, 262: 125234. 
	  
	  \bibitem{ref16} 王雅博, 诸凯, 崔卓, 等. 冷却水进出方式对芯片散热器换热性能影响[J]. 机械工程学报, 2018, 54(10): 188-194.
	  
	  \bibitem{ref17} Choi W W, Bang S H, Shin H C, et al. Experimental and numerical study of a high-performance air-cooled N-type manifold microchannel heat sink for large-area applications[J]. International Communications in Heat and Mass Transfer, 2026, 175: 110986.

	 
	  \bibitem{ref20} Indhe M J, Sur A, Yadav A K, et al. Integrating CFD, molecular dynamics, and AI techniques for thermal management of heat transfer devices in microelectronic chip cooling using nano-enhanced PCMs and hybrid cooling systems[J]. International Journal of Thermofluids, 2026, 32: 101587. 
	  \bibitem{ref21} Nie C, Wang H, Li H, et al. Flow and heat transfer characteristics of pin-fin microchannel heat sink with associated diversion structure[J]. International Communications in Heat and Mass Transfer, 2026, 175: 111004.
	  \bibitem{yin} Yin Y, Wang Z, Ma S, et al. A novel coaxial open annular TSV pin-fin embedded in microchannels to comprehensively improve the thermal stability, signal integrity, and heat transfer of 3D integrated circuits[J]. International Communications in Heat and Mass Transfer, 2026, 172: 110619. 
	  \bibitem{seong} Seong H B, Gwon J G, Lee M S, et al. Numerical study and ANN analysis on heat transfer and flow characteristics of microchannel heat sink with inclined rib–pin fin in in-line and staggered arrangements[J]. International Communications in Heat and Mass Transfer, 2026, 178: 111593. 
	  \bibitem{ref24} Lorenzini D, Green C, Sarvey T E, et al. Embedded single phase microfluidic thermal management for non-uniform heating and hotspots using microgaps with variable pin fin clustering[J]. International Journal of Heat and Mass Transfer, 2016, 103: 1359-1370. 
	 
	  \bibitem{yuan} 袁永帅, 等. 当前芯片散热技术方法和特点研讨[J]. 2024.
	  \bibitem{xiong} 熊园园, 等. 基于内嵌式微通道芯片散热结构设计研究综述[J]. 2024.
	\end{thebibliography}
	\textcolor{red}{昊洋借助AI和申报书适当补充文献，并让AI辅助导入新文献，文献总数在10篇左右}
\end{document}
\bibitem{xiong} 熊园园, 等. 基于内嵌式微通道芯片散热结构设计研究综述[J]. 2024.
\end{thebibliography}
\textcolor{red}{昊洋借助AI和申报书适当补充文献，并让AI辅助导入新文献，文献总数在10篇左右}
\end{document}
